% !TEX root = main.tex
%======================================================================
%  M1 Scoping note --- Positioning / relation to prior work + novelty
%======================================================================
\section{Positioning: why this is a genuine gap}
\label{sec:positioning}

The setting of Section~\ref{sec:problem} --- \emph{model-based} DFO with
fully-linear models of the objective \emph{and} of \emph{relaxable inequality}
constraints, aiming at KKT stationarity under MFCQ --- sits in a gap left open by
the closest lines of work. We locate each precedent precisely and identify what
it does not cover. Every claim below is stated against the source PDF in
\path{articles/}.

\paragraph{Trust-funnel lineage (derivative-based).}
\citet{GouldToint2010} introduced the trust-funnel method: equality-constrained
nonlinear programming ``without a penalty function or a filter,'' using a
normal/tangential step decomposition and a monotone infeasibility cap (the
funnel). We adopt this normal/tangential + funnel machinery
(Section~\ref{sec:funnel}) but replace derivatives with fully-linear models and
target inequalities.

\paragraph{Derivative-free trust-funnel (equalities).}
\citet{SampaioToint2015} is the primary precedent: a derivative-free trust-funnel
method that builds fully-linear models and proves first-order convergence --- but
for \emph{equality} constraints $c(x)=0$. The companion numerical study
\citep{SampaioToint2016} reports experience extending the approach to problems
``with both equality and inequality constraints and \dots\ simple bounds,''
handling inequalities by the usual device of adding slack variables and using
``neither filter nor penalty function.'' Our project differs in two ways that
matter for the theory: (i) we treat inequalities \emph{directly} through the
violation measure $\theta=\|[c]_+\|$ rather than by slack-variable
equality reformulation, so the model geometry lives in the original $x$-space that
the simulation samples; and (ii) our target is a \emph{proof} of KKT convergence
under MFCQ for that direct inequality treatment, not a numerical extension.

\paragraph{Model-based DFO with feasible inequality iterates (the contrast).}
NOWPAC \citep{AugustinMarzouk2017} and its stochastic generalization (S)NOWPAC
\citep{Menhorn2022} build fully-linear models of the objective and of
\emph{inequality} constraints inside a trust region --- superficially the closest
match. The decisive difference is that NOWPAC enforces \emph{feasible iterates}:
it adds a convex offset $h_k$ to the constraints, the ``inner boundary path,''
which keeps iterates strictly inside the feasible region and is precisely what
their convergence argument relies on \citep[\S4]{Menhorn2022}. That design
targets \emph{unrelaxable} constraints (points outside $\feas$ may not be usable)
--- the opposite of our Assumption~\ref{ass:relaxable}. Requiring feasibility
throughout forbids the infeasible model-sample points and normal-step excursions
that our funnel both permits and needs. So NOWPAC solves an adjacent but distinct
problem, and its proof technique does not transfer to the relaxable setting.

\paragraph{Feasibility-restoration alternatives.}
Inexact-restoration DFO \citep{Birgin2021} handles constraints by alternating a
feasibility-restoration phase with an optimization phase. It is a legitimate
alternative architecture for relaxable constraints and a useful contrast, but it
is not a trust-funnel method and does not build a single fully-linear model set
used jointly by a normal and a tangential step; we mention it as related rather
than as the framework we extend.

\paragraph{Constraint taxonomy and the DFO backdrop.}
\citet{LeDigabelWild2015} give the taxonomy (quantifiable/relaxable/simulation)
that names our constraint class and explains why relaxability is the enabling
property. \citet{LMW2019} survey constrained DFO and the fully-linear-model /
trust-region machinery, and \citet[Ch.~6,10,11]{CSV2009} provide the
$\Lambda$-poisedness geometry (uniform fully-linearity constants) and the
unconstrained trust-region convergence template that our criticality step and
$\kappa_{\mathrm{eg}}\Delta_k\to0$ argument specialize.

\subsection{Novelty statement}
\label{sec:novelty}

To our knowledge, no existing method provides a \emph{model-based} trust-funnel
algorithm --- with fully-linear models of \emph{both} the objective and the
\emph{inequality} constraints --- for \emph{relaxable} simulation constraints,
together with a proof of convergence to a KKT stationary point under MFCQ, and
\emph{without} a penalty parameter or a feasible-iterate requirement.

\medskip\noindent
Concretely, the contribution combines
\begin{enumerate}
  \item the DFO trust-funnel of \citet{SampaioToint2015} --- but extended from
        \emph{equalities} to \emph{inequalities} via the violation measure
        $\theta=\|[c]_+\|$ handled \emph{directly} in $x$-space;
  \item the fully-linear objective-and-inequality-constraint modeling of NOWPAC
        \citep{AugustinMarzouk2017,Menhorn2022} --- but \emph{without} its
        feasible-iterate inner-boundary requirement, using relaxability
        (Assumption~\ref{ass:relaxable}) to sample and step through infeasible
        points; and
  \item a KKT-under-MFCQ convergence guarantee, obtained without any penalty
        parameter (contrast the penalty $h$-functions in the repository, e.g.\
        \path{general_h_funs.py}, which this project deliberately does \emph{not}
        use).
\end{enumerate}

\subsection{Summary table}
\label{sec:table}

\begin{center}\small
\begin{tabular}{@{}lccccc@{}}
\toprule
Method & DF? & FL model of $c$ & Constraint type & Iterates & Convergence \\
\midrule
Gould--Toint \citeyearpar{GouldToint2010}   & no  & --- & equality            & funnel  & 1st-order \\
Sampaio--Toint \citeyearpar{SampaioToint2015} & yes & yes & equality           & funnel  & 1st-order \\
NOWPAC / (S)NOWPAC                            & yes & yes & inequality          & feasible-only & yes \\
Inexact restoration \citeyearpar{Birgin2021} & yes & (rest.) & general         & restoration & yes \\
\textbf{This project}                        & \textbf{yes} & \textbf{yes} & \textbf{relaxable ineq.} & \textbf{funnel (infeasible OK)} & \textbf{KKT/MFCQ (target)} \\
\bottomrule
\end{tabular}
\end{center}

\noindent
``FL model of $c$'' = fully-linear model of the constraints; ``DF'' =
derivative-free. The last row is the deliverable of milestones M2--M3.
